% %%%%%%%%%%%%%%%%%%%%%%%%%%%%%%%%%%%%%%%%%%%%%%%%%%%%%%%%%%%%%%%%%%%%%
% %% SpDCache %%
% Conclusion
% %%%%%%%%%%%%%%%%%%%%%%%%%%%%%%%%%%%%%%%%%%%%%%%%%%%%%%%%%%%%%%%%%%%%%

\section{Conclusion}
\label{sec:spdc:conclu}

% \lettrine{M}{odern}

We have presented the concept of SpDCache, a \emph{reduction cache} that has optimized storage and compute techniques for the sparser and denser regions of \emph{banded} matrices.
This approach does not reduce the number of compute operations, however, it significantly decreases memory traffic ($8.1\times$).
These benefits can be attributed to three techniques: (1) blocking the dense part of the \emph{band} in the \emph{IBRC}, (2) using fine-grained storage in the \emph{OBRT} and (3) off-loading the \emph{reductions} for isolated elements to the downstream caches, to prevent local \emph{evictions} and reduce traffic.

Exploiting the coarse-grain structure of matrices is key to design hardware optimized for sparse matrix kernels.
In future work, we plan to also take advantage of local fine-grain density in the \emph{OBR}, to further reduce memory traffic, and evaluate speedup on a clock-accurate model.

\bigskip\bigskip

\begin{highlight}{Takeaways}{Mulberry}
    \raggedright
    \colorbf{Mulberry}{ra} tata
\end{highlight}

\clearpage